% ===== Depression: "Dealing with the black dog" (~20 min) =====
\topicopener{raincloud}{Դեպրեսիա}{Ինչպես ապրել սև շան հետ}{Դեպրեսիան և դրա շուրջ թյուր կարծիքները}{sad}

\begin{frame}{Մի բառ՝ երկու բան}
  \begin{columns}[T]
    \begin{column}{0.47\textwidth}
      \begin{card}[colback=butter!45!cream]
        \emoji[1.3em]{sunbehindcloud}\ {\blackfont «Դեպրեսիա ունեմ»}\par\smallskip
        Առօրյա խոսքում՝ «մի քիչ տխուր եմ»։ Անցնում է
      \end{card}
    \end{column}
    \begin{column}{0.47\textwidth}
      \onslide<2->{%
      \begin{card}[colback=lavender!45!cardwhite]
        \emoji[1.3em]{raincloud}\ {\blackfont Կլինիկական դեպրեսիա}\par\smallskip
        Իրական, ծանր տրամադրության խանգարում (major depressive disorder)
      \end{card}}
    \end{column}
  \end{columns}
  \vspace{3mm}
  \onslide<3->{%
  \begin{bookquote}
    Դեպրեսիա ունեցողին լուրջ չընդունել, որովհետև դու էլ ես տխրել ու անցել է, նույնն է, թե ձեռքն անդամահատված մարդուն ասես. «Ես էլ եմ մի անգամ թղթով մատս կտրել»։
  \end{bookquote}}
\end{frame}

\begin{frame}{Ինչու ենք թերագնահատում}
  \begin{columns}[c]
    \begin{column}{0.56\textwidth}
      \begin{itemize}
        \item Դատում ենք \textbf{մեր} փորձով. «բոլորս էլ ենք տխրում, հետո անցնում է»
        \item Ուղեղը մեր փորձը մեծացնում է, իսկ ուրիշներինը՝ փոքրացնում։ Հիշո՞ւմ ես եսասեր հիշողությունը առաջին վեբինարից
        \item<2-> Իսկ իրական դեպրեսիան կարող է այնքան ծանր լինել, որ մարդը սկսի մտածել կյանքից հեռանալու մասին
      \end{itemize}
    \end{column}
    \begin{column}{0.38\textwidth}
      \onslide<3->{%
      \begin{card}[colback=mint!35!cardwhite]
        \centering\bigemoji[1cm]{heart}\par\smallskip
        \raggedright Դեպրեսիան իրական, լուրջ վիճակ է, որն արժանի է \textbf{կարեկցանքի ու հարգանքի}, ոչ թե արհամարհանքի
      \end{card}}
    \end{column}
  \end{columns}
\end{frame}

\begin{frame}{«Եսասեր», «ծույլ»}
  \relax{\small Դեպրեսիա ունեցողներին ու ինքնասպանությամբ մահացածներին հաճախ այսպես են անվանում։ Ինչո՞ւ}\par\smallskip
  \centering
  \begin{tikzpicture}[pbox/.style={flow=3.5cm, font=\small}, yscale=0.75]
    \node[pbox, fill=butter!50!cardwhite] (s) at (0,0) {Ուրիշի ծանր\\ դեպրեսիան};
    \node[pbox, fill=lavender!40!cardwhite, visible on=<2->] (a1) at (4.4,0.75) {«Իրական հիվանդություն է»};
    \node[pbox, fill=brainpink!35!cardwhite, visible on=<2->] (a2) at (8.6,0.75) {\emoji{fearful}\ Ուրեմն ինձ հետ էլ կարող է պատահել};
    \node[pbox, fill=lavender!40!cardwhite, visible on=<3->] (b1) at (4.4,-0.75) {«Նա ուղղակի եսասեր է»};
    \node[pbox, fill=mint!40!cardwhite, visible on=<3->] (b2) at (8.6,-0.75) {\emoji{slightsmile}\ Ինձ հետ չի պատահի};
    \draw[arr, visible on=<2->] (s.east) -- (a1.west);
    \draw[arr, visible on=<2->] (a1) -- (a2);
    \draw[arr, visible on=<3->] (s.east) -- (b1.west);
    \draw[arr, visible on=<3->] (b1) -- (b2);
  \end{tikzpicture}\par\smallskip
  \raggedright
  \onslide<4->{%
  \begin{bookquote}
    Սա մեկ բացատրություն է։ Մյուսն այն է, որ որոշ մարդիկ ուղղակի տգետ տխմարներ են։
  \end{bookquote}}
\end{frame}

\begin{frame}{Ռոբին Ուիլյամսի մահից հետո}
  \begin{columns}[T]
    \begin{column}{0.5\textwidth}
      \begin{itemize}
        \item 2014. ինքնասպանությամբ մահացավ սիրված դերասան ու կատակերգու Ռոբին Ուիլյամսը \emoji{candle}
        \item Ջերմ խոսքերի կողքին՝ «Եսասիրություն է», «Ամեն ինչ ուներ...», նաև հայտնի մարդկանցից ու լրատվամիջոցներից
      \end{itemize}
      \onslide<2->{%
      \begin{bookquote}
        Եթե դու էլ ես այդպես ասել, ներիր, բայց սխալվում ես։
      \end{bookquote}}
    \end{column}
    \begin{column}{0.44\textwidth}
      \onslide<3->{%
      \begin{scienceupdate}
        Հետմահու հետազոտությունը ցույց տվեց, որ նա ուներ նաև \textbf{Լևիի մարմնիկներով դեմենցիա}. այն կարող է անհանգստություն, պարանոյա ու դեպրեսիա առաջացնել։\par\smallskip
        Դրսից չենք տեսնում, թե ինչ է կատարվում ուրիշի ուղեղում։
      \end{scienceupdate}}
    \end{column}
  \end{columns}
  \only<3->{\source{S. Schneider Williams, Neurology, 2016}}
\end{frame}

\begin{frame}{Դեպրեսիան տարբեր դեմքեր ունի}
  \begin{columns}[T]
    \begin{column}{0.47\textwidth}
      \begin{card}[colback=lavender!35!cardwhite]\emoji[1.3em]{pensive}\ \textbf{Անհույս թախիծ}, որից չես կարողանում դուրս գալ\end{card}
      \onslide<2->{\begin{card}[colback=peach!35!cardwhite]\emoji[1.3em]{fearful}\ \textbf{Ուժեղ անհանգստություն}, մոտալուտ աղետի զգացում\end{card}}
    \end{column}
    \begin{column}{0.47\textwidth}
      \onslide<3->{\begin{card}[colback=butter!45!cream]\emoji[1.3em]{neutral}\ \textbf{Դատարկություն}. ոչ մի զգացմունք, ինչ էլ որ կատարվի\end{card}}
      \onslide<4->{\begin{card}[colback=brainpink!25!cardwhite]\emoji[1.3em]{angry}\ \textbf{Մշտական զայրույթ} ու անհանգստություն (ավելի հաճախ՝ տղամարդկանց մոտ)\end{card}}
    \end{column}
  \end{columns}
  \vspace{3mm}
  \onslide<5->{\centering Դրա համար էլ դեպրեսիայի \textbf{մեկ պատճառ} գտնելը շատ դժվար է։\par}
\end{frame}

\begin{frame}{Մոնոամինային վարկածը}
  \begin{columns}[c]
    \begin{column}{0.42\textwidth}
      \begin{itemize}
        \item Շատ նեյրոմեդիատորներ \textbf{մոնոամիններ} են. սերոտոնին, դոֆամին, նորադրենալին
        \item Դեպրեսիայի ժամանակ դրանք կարծես \textbf{քիչ} են
        \item Հակադեպրեսանտների մեծ մասն ավելացնում է դրանք
        \item<2-> Ամենատարածվածը՝ \textbf{SSRI}-ները (selective serotonin reuptake inhibitors)
      \end{itemize}
    \end{column}
    \begin{column}{0.56\textwidth}
      \centering
      % A serotonin synapse: the reuptake "pump" pulls serotonin back; an SSRI blocks it (slide 2).
% Schematic, after the book's description of SSRIs. Use inside a frame (overlays: <2->).
\begin{tikzpicture}[scale=1.12,
    lbl/.style={font=\scriptsize, text=plum, align=center, inner sep=1pt},
  ]
  % sending and receiving neurons (straight membranes face the gap)
  \filldraw[fill=brainpink!70!cardwhite, draw=plum, line width=1pt, rounded corners=5mm] (-3.3,-1.35) rectangle (-0.45,1.35);
  \filldraw[fill=lavender!80!cardwhite, draw=plum, line width=1pt, rounded corners=5mm] (0.45,-1.35) rectangle (3.3,1.35);
  \node[lbl, font=\scriptsize\bfseries] at (-2.15,1.0) {նեյրոն 1};
  \node[lbl, font=\scriptsize\bfseries] at (2.15,1.0) {նեյրոն 2};
  % vesicles full of serotonin
  \foreach \x/\y in {-2.3/0.15, -1.45/0.55, -1.7/-0.5}{
    \filldraw[fill=cardwhite, draw=plum, line width=0.6pt] (\x,\y) circle (0.3);
    \foreach \dx/\dy in {-0.1/0.06, 0.08/0.1, 0.0/-0.1}{\fill[amBlue] (\x+\dx,\y+\dy) circle (0.055);}
  }
  % receptors on neuron 2
  \foreach \y in {0.55, 0.0, -0.55}{
    \draw[amOrange!90!black, line width=1.6pt] (0.45,\y+0.15) -- (0.62,\y+0.15) -- (0.62,\y-0.15) -- (0.45,\y-0.15);}
  % serotonin in the gap
  \foreach \x/\y in {-0.15/0.85, 0.12/0.42, -0.2/0.1, 0.2/-0.25, -0.05/-0.55}{\fill[amBlue] (\x,\y) circle (0.065);}
  \foreach \x/\y in {0.15/0.7, -0.22/0.55, 0.05/0.15, -0.12/-0.3, 0.22/0.05, -0.25/-0.1, 0.1/-0.75}{\fill[amBlue, visible on=<2->] (\x,\y) circle (0.065);}
  % reuptake pump on neuron 1
  \filldraw[fill=mint, draw=plum, line width=0.8pt, rounded corners=0.6mm] (-0.6,-1.15) rectangle (-0.3,-0.82);
  \draw[arr=mintdark, line width=1.1pt, visible on=<1>] (-0.05,-0.98) .. controls (-0.3,-0.98) and (-0.6,-0.98) .. (-1.0,-0.98);
  % the SSRI blocks the pump
  \filldraw[visible on=<2->, fill=amRed, draw=plum, line width=0.7pt] (-0.45,-0.98) circle (2.2mm);
  \draw[visible on=<2->, white, line width=1.3pt, line cap=round] (-0.55,-1.08) -- (-0.35,-0.88) (-0.55,-0.88) -- (-0.35,-1.08);
  % labels
  \node[lbl, anchor=south] (s) at (0,1.55) {\textcolor{amBlue}{\textbf{սերոտոնին}}};
  \draw[plum!60, line width=0.5pt] (s.south) -- (-0.12,0.95);
  \node[lbl] (r) at (2.05,-0.2) {\textcolor{amOrange!80!black}{\textbf{ընկալիչներ}}};
  \draw[plum!60, line width=0.5pt] (r.west) -- (0.68,0.0);
  \node[lbl, anchor=north, visible on=<1>] at (-0.45,-1.45) {հետ է «քաշվում»\\ (reuptake)};
  \node[lbl, anchor=north, visible on=<2->, text=amRed] at (-0.45,-1.45) {\textbf{SSRI-ն փակում է «պոմպը»}\\ սերոտոնինը ճեղքում մնում է երկար};
\end{tikzpicture}

    \end{column}
  \end{columns}
\end{frame}

\begin{frame}{Ինչու է այս վարկածը քննադատվում}
  \begin{columns}[c]
    \begin{column}{0.46\textwidth}
      \centering
      % Schematic, not data: an SSRI raises serotonin within a day, but mood improves only after weeks
% (the book: "like filling your car's empty tank and it working again only a month later").
\begin{tikzpicture}
  \begin{axis}[
      width=6.4cm, height=4.4cm,
      xmin=0, xmax=6.3, ymin=0, ymax=1.12,
      xtick={0,1,2,3,4,5,6}, ytick=\empty,
      xlabel={շաբաթներ դեղն սկսելուց հետո}, xlabel style={font=\scriptsize, text=plum},
      tick label style={font=\scriptsize, text=plum}, axis line style={plum},
      axis x line*=bottom, axis y line*=left, clip=false,
    ]
    \addplot[amBlue, line width=1.6pt] coordinates {(0,0.15) (0.12,0.88) (6,0.88)};
    \node[anchor=south west, font=\blackfont\scriptsize, text=amBlue, inner sep=1pt] at (axis cs:0.25,0.9) {սերոտոնին՝ միանգամից};
    \addplot[amOrange, line width=1.6pt, domain=0:6, samples=60, visible on=<2->] {0.15 + 0.6/(1+exp(-2.2*(x-3.2)))};
    \node[anchor=north west, font=\blackfont\scriptsize, text=amOrange!85!black, inner sep=1pt, visible on=<2->,
          align=left] at (axis cs:3.6,0.55) {տրամադրություն՝\\ շաբաթներ անց};
  \end{axis}
\end{tikzpicture}

    \end{column}
    \begin{column}{0.5\textwidth}
      \begin{itemize}
        \item «Ավելի շատ կանաչ է պետք»՝ ասել հին նկար վերականգնողին։ Գուցե ճիշտ է, բայց \textbf{ինչ անել}՝ չի ասում \emoji{palette}
        \item<2-> Ասես բաքը լցրիր, իսկ մեքենան միացավ \textbf{մեկ ամիս} անց \emoji{car}
        \item<3-> «Փչացած» համակարգ չի գտնվել. կան նաև հակադեպրեսանտներ, որ մոնոամիններին չեն դիպչում
      \end{itemize}
    \end{column}
  \end{columns}
\end{frame}

\begin{frame}{«Քիմիական անհավասարակշռությու՞ն»}
  \begin{columns}[c]
    \begin{column}{0.74\textwidth}
      \begin{scienceupdate}
        Բըրնեթն արդեն 2016-ին գրում էր, որ դեպրեսիան «քիմիական անհավասարակշռությունից» ավելին է։\par\smallskip
        2022-ի մեծ ակնարկը համոզիչ ապացույց չգտավ, որ դեպրեսիան առաջանում է սերոտոնինի \textbf{պակասից}։ Ակնարկը վիճարկվեց, բայց շատերը համաձայն են. «սերոտոնինի պակասը» չափազանց պարզ բացատրություն է։\par\smallskip
        \onslide<2->{Սա \textbf{չի նշանակում}, որ հակադեպրեսանտները չեն աշխատում. մեծ վերլուծություններում դրանք պլացեբոյից ավելի արդյունավետ են։}
      \end{scienceupdate}
    \end{column}
    \begin{column}{0.22\textwidth}
      \centering\mascot[2.3cm]{thinking}
    \end{column}
  \end{columns}
  \source{Moncrieff et al., 2022; Jauhar et al., 2023; Cipriani et al., Lancet, 2018}
\end{frame}

\begin{frame}{Այլ վարկածներ}
  \begin{columns}[c]
    \begin{column}{0.44\textwidth}
      \centering
      % Medial view of the brain (front = left; images/brain_medial.png, P. J. Lynch, Wikimedia Commons, CC BY 2.5)
% with the areas the depression theories point at. Positions are fractions of the image (0,0 = bottom left).
\begin{tikzpicture}
  \node[anchor=south west, inner sep=0] (b) at (0,0) {\includegraphics[width=5cm]{brain_medial.png}};
  \begin{scope}[x={(b.south east)}, y={(b.north west)}]
    \node[braindot=amRed] (acc) at (0.30,0.715) {};
    \node[brainlbl=brainpink!40!cardwhite, anchor=south] (al) at (0.42,1.02) {առաջային գոտկային կեղև};
    \draw[plum, line width=0.6pt] (acc) -- (al.south);
    \node[braindot=amBlue, visible on=<2->] (hyp) at (0.435,0.415) {};
    \node[braindot=amBlue, visible on=<2->] (pit) at (0.355,0.265) {};
    \node[brainlbl=lavender!50!cardwhite, anchor=north, visible on=<2->] (hl) at (0.36,-0.03) {հիպոթալամուս ու հիպոֆիզ (HPA)};
    \draw[plum, line width=0.6pt, visible on=<2->] (hyp) -- (hl.north);
    \draw[plum, line width=0.6pt, visible on=<2->] (pit) -- (hl.north);
  \end{scope}
\end{tikzpicture}

    \end{column}
    \begin{column}{0.52\textwidth}
      \begin{card}[colback=lavender!35!cardwhite]\emoji{moon}\ \textbf{Քուն}. սերոտոնինը կարգավորում է ներքին ժամացույցը, իսկ դեպրեսիան խախտում է քունը\end{card}
      \begin{card}[colback=brainpink!25!cardwhite]\textcolor{amRed}{\Large$\bullet$}\ \textbf{Առաջային գոտկային կեղև} (anterior cingulate)՝ ուղեղի «շվեյցարական դանակը»։ Դեպրեսիայի ժամանակ ավելի ակտիվ է\end{card}
      \onslide<2->{\begin{card}[colback=mint!30!cardwhite]\textcolor{amBlue}{\Large$\bullet$}\ \textbf{Սթրեսի առանցքը} (HPA axis)\end{card}}
    \end{column}
  \end{columns}
\end{frame}

\begin{frame}{Ուղեղը, որ կորցրել է ճկունությունը}
  \begin{columns}[c]
    \begin{column}{0.54\textwidth}
      \begin{itemize}
        \item \textbf{Նեյրոպլաստիկություն}՝ նեյրոնների միջև նոր կապեր ստեղծելու ունակություն։ Դրանով ենք սովորում ու հարմարվում
        \item Դեպրեսիայի ժամանակ այն \textbf{թուլանում} է. ուղեղը «կարծրանում» է, ինչպես շատ երկար դրսում մնացած կարկանդակ \emoji{cake}
        \item Վատ բանից հետո դժվար է առաջ շարժվել, ու դեպրեսիան մնում է
      \end{itemize}
    \end{column}
    \begin{column}{0.4\textwidth}
      \onslide<2->{%
      \begin{card}[colback=mint!35!cardwhite]
        \centering\bigemoji[1cm]{seedling}\par\smallskip
        \raggedright Հակադեպրեսանտները հաճախ \textbf{ավելացնում} են պլաստիկությունը։ Գուցե դրա համար է օգուտը շաբաթներ անց գալիս։\par\smallskip
        Դա վառելիք լցնել չէ, այլ բույսը \textbf{պարարտացնել}
      \end{card}}
    \end{column}
  \end{columns}
\end{frame}

\begin{frame}{Պատճա՞ռ, թե՞ հետևանք}
  \begin{columns}[c]
    \begin{column}{0.56\textwidth}
      \begin{itemize}
        \item Այս բոլոր վարկածները կարող են լինել դեպրեսիայի և՛ պատճառ, և՛ հետևանք։ Հետազոտությունները շարունակվում են
        \item<2-> Դեպրեսիան վատացնում է նաև \textbf{մտածողությունը}. թեստում «չեմ հիշում»-ն ու «ուժ չունեմ փորձելու»-ն նույն արդյունքն են տալիս
        \item<3-> Իսկ դեմենցիայի ախտորոշումը հաճախ \textbf{դեպրեսիա} է բերում
      \end{itemize}
    \end{column}
    \begin{column}{0.4\textwidth}
      \centering
      \onslide<2->{%
      \begin{tikzpicture}
        \node (t) at (0,0) {\bigemoji[1.2cm]{clipboard}};
        \node[font=\blackfont\small, text=plum] at (0,-0.95) {թեստ՝ ցածր արդյունք};
        \node[flow=2.3cm, fill=lavender!40!cardwhite] (a) at (-1.45,-2.1) {դեմենցիա՞};
        \node[flow=2.3cm, fill=peach!45!cardwhite] (b) at (1.45,-2.1) {դեպրեսիա՞};
        \draw[arr] (-0.3,-1.25) -- (a.north);
        \draw[arr] (0.3,-1.25) -- (b.north);
      \end{tikzpicture}\par\medskip
      {\small Բժիշկներին սովորեցնում են տարբերել. բուժումը շատ տարբեր է}}
    \end{column}
  \end{columns}
\end{frame}

\begin{frame}{Աչքը որսում է վատը}
  \relax{\small Դեպրեսիա ունեցողները բառերի ցուցակում ավելի շատ ուշադրություն են դարձնում \textbf{տհաճ} բառերին։}\par\medskip
  \centering
  % The same words, two filters: the usual ego-protecting bias vs depression's negative bias (book: Epp et al., 2012).
% Use inside a frame (overlays: <2->).
\begin{tikzpicture}[
    w/.style={draw=plum, line width=0.7pt, fill=cardwhite, rounded corners=1.5mm, font=\small, text=plum, inner sep=3pt, minimum width=2.3cm},
  ]
  \node[w] (w1) at (0,1.05) {խոտ};
  \node[w] (w2) at (0,0.35) {սպանություն};
  \node[w] (w3) at (0,-0.35) {գովեստ};
  \node[w] (w4) at (0,-1.05) {ձախողում};
  \node[font=\scriptsize, text=plum!70!cream] at (0,1.65) {ինչ է կատարվում};
  % usual brain
  \node[flow=4.6cm, fill=mint!40!cardwhite, anchor=west] (n) at (2.4,0.95)
    {\textbf{Սովորաբար}\\ {\color{mintdark}\textbf{գովեստ}}, խոտ\\ {\color{plum!40!cream} սպանություն, ձախողում}};
  % depression
  \node[flow=4.6cm, fill=brainpink!35!cardwhite, anchor=west, visible on=<2->] (d) at (2.4,-0.95)
    {\textbf{Դեպրեսիայի ժամանակ}\\ {\color{pinkdark}\textbf{սպանություն, ձախողում}}\\ {\color{plum!40!cream} գովեստ, խոտ}};
  \draw[arr] (1.3,0.3) -- (n.west);
  \draw[arr, visible on=<2->] (1.3,-0.3) -- (d.west);
\end{tikzpicture}
\par\medskip
  \onslide<3->{\small Սովորական եսակենտրոն շեղումը շրջվում է, ու դեպրեսիայից դուրս գալը շատ դժվար է դառնում։\par}
\end{frame}

\begin{frame}{Վայրընթաց պարույր}
  \begin{columns}[c]
    \begin{column}{0.58\textwidth}
      % Burnett's "downward spiral": problems pile up one after another. Steps reveal on slides 1-5.
\begin{tikzpicture}[st/.style={flow=5.6cm, align=left, inner sep=3.5pt, font=\small}]
  \node[st, fill=butter!55!cardwhite] (a) at (0,2.0) {\emoji{boss}\ Կորցնում ես աշխատանքդ};
  \node[st, fill=peach!50!cardwhite, visible on=<2->] (b) at (0.6,1.0) {\emoji{brokenheart}\ Զուգընկերդ հեռանում է};
  \node[st, fill=peach!70!cardwhite, visible on=<3->] (c) at (1.2,0) {\emoji{candle}\ Հարազատդ մահանում է};
  \node[st, fill=brainpink!45!cardwhite, visible on=<4->] (d) at (1.8,-1.0) {\emoji{warning}\ Թաղումից գալիս թալանում են};
  \node[st, fill=brainpink!70!cardwhite, visible on=<5->] (e) at (2.4,-2.0) {\emoji{wine}\ Մեկուսացում, ալկոհոլ...};
  \draw[arr, visible on=<2->] ([xshift=3mm]a.south west) |- (b.west);
  \draw[arr, visible on=<3->] ([xshift=3mm]b.south west) |- (c.west);
  \draw[arr, visible on=<4->] ([xshift=3mm]c.south west) |- (d.west);
  \draw[arr, visible on=<5->] ([xshift=3mm]d.south west) |- (e.west);
\end{tikzpicture}

    \end{column}
    \begin{column}{0.38\textwidth}
      \onslide<6->{%
      \begin{card}[colback=lavender!35!cardwhite]
        Ուղեղի հարմար համոզմունքները («աշխարհն արդար է», «ինձ հետ վատ բան չի լինի») \textbf{փշրվում} են։ Վերահսկողություն չկա։\par\smallskip
        Ալկոհոլն ու մեկուսացումը մի պահ թեթևացնում են, բայց ավելի են ծանրաբեռնում ուղեղը
      \end{card}}
    \end{column}
  \end{columns}
\end{frame}

\begin{frame}{Ռիսկի գործոններ. կայծակի օրինակը}
  \begin{columns}[c]
    \begin{column}{0.44\textwidth}
      \centering
      \begin{tikzpicture}
        \node[flow=2.4cm, fill=mint!35!cardwhite] (in) at (-1.4,0) {\bigemoji[0.9cm]{house}\\ տանը};
        \node[flow=2.4cm, fill=brainpink!30!cardwhite] (out) at (1.4,0) {\bigemoji[0.9cm]{lightning}\bigemoji[0.9cm]{tree}\\ դրսում, ծառի տակ};
      \end{tikzpicture}\par\medskip
      {\small Ռիսկը \textbf{ավելացնում} է հավանականությունը, բայց ոչինչ չի երաշխավորում}
    \end{column}
    \begin{column}{0.52\textwidth}
      \begin{itemize}
        \item Աղքատություն, վտանգավոր թաղամաս, մենակություն
        \item<2-> Ծխելը մեծացնում է թոքի քաղցկեղի ռիսկը, բայց այն լինում է նաև չծխողների մոտ
        \item<3-> Անհատականություն ու \textbf{գեներ}. դեպրեսիան մոտ 37\%-ով ժառանգական է
      \end{itemize}
    \end{column}
  \end{columns}
  \only<3->{\source{Sullivan, Neale \& Kendler, Am. J. Psychiatry, 2000}}
\end{frame}

\begin{frame}{Առասպել, թե՞ փաստ}
  \mythorfact{Հարուստ ու հայտնի մարդը դեպրեսիա չի ունենա}{ԱՌԱՍՊԵԼ}{pinkdark}{%
    \begin{card}[colback=brainpink!25!cardwhite]
      Հաջողությունը պատվաստ չէ։ Կինոյում նկարահանվելուց ուղեղի աշխատանքը չի փոխվում։
    \end{card}}
\end{frame}

\begin{frame}{«Բայց չէ՞ որ ամեն ինչ ունի»}
  \begin{columns}[c]
    \begin{column}{0.54\textwidth}
      \begin{itemize}
        \item Դեպրեսիայի դեմ պայքարը կարող է \textbf{մղել} դեպի հաջողություն. «ծաղրածուի արցունքներ», շատ կատակերգուներ, Վան Գոգ \emoji{palette}
        \item<2-> Նպատակին հասնելը կարող է խլել կյանքի \textbf{իմաստը}
        \item<3-> Գուցե ճանապարհին կորցրել ես սիրելի մարդկանց
      \end{itemize}
    \end{column}
    \begin{column}{0.4\textwidth}
      \onslide<4->{%
      \begin{bookquote}
        Եթե չես կարող դեպրեսիա ունենալ, որովհետև ուրիշներից լավ ես ապրում, ուրեմն միայն Երկրի ամենադժբախտ մարդը պետք է դեպրեսիա ունենա։
      \end{bookquote}}
    \end{column}
  \end{columns}
\end{frame}

\begin{frame}{Դեպրեսիան տրամաբանական չէ}
  \begin{columns}[c]
    \begin{column}{0.52\textwidth}
      \begin{itemize}
        \item «Եսասեր» ասողները կարծես պատկերացնում են, որ մարդը «կողմ ու դեմ» աղյուսակ է կազմում
        \item Բայց դեպրեսիան փոխում է հենց \textbf{մտածողությունը}. ամեն ինչ անցնում է ուղեղի «ֆիլտրով»
        \item<3-> Այն կարող է մարդուն համոզել, որ նա \textbf{բեռ} է ուրիշների համար։ Սա հիվանդությունն է խոսում, ոչ թե սառը հաշվարկը
      \end{itemize}
    \end{column}
    \begin{column}{0.44\textwidth}
      \onslide<2->{%
      \begin{card}[colback=lavender!35!cardwhite]\emoji[1.3em]{wave}\ Խեղդվողը չի «շնչում օդը» այնպես, ինչպես ափին կանգնածը\end{card}}
      \onslide<4->{%
      \begin{card}[colback=peach!35!cardwhite]\emoji[1.3em]{bandage}\ Դեպրեսիա ունեցողից տրամաբանություն պահանջելը նույնն է, թե կոտրված ոտքով մարդուն ասես՝ «նորմալ քայլիր»\end{card}}
    \end{column}
  \end{columns}
\end{frame}

\begin{frame}{«Ընտրությո՞ւն» է}
  \begin{columns}[c]
    \begin{column}{0.54\textwidth}
      \begin{itemize}
        \item Մեղադրանքները ենթադրում են, թե մարդը \textbf{ընտրել} է իր վիճակը։ Ինչո՞ւ ու ինչպե՞ս, ոչ ոք չի բացատրում
        \item Տառապանքը, որ հաղթում է միլիոնավոր տարիների գոյատևման բնազդին, «հեշտ ելք» չես անվանի
        \item<2-> Ուղեղը սիրում է \textbf{հետևողականություն}, ու կարծիքը փոխելը դժվար է։ Իսկ ուրիշին մեղադրելը՝ ինքդ քեզ լավ զգալու համար, եսասիրության շատ ավելի լավ օրինակ է
      \end{itemize}
    \end{column}
    \begin{column}{0.4\textwidth}
      \onslide<3->{%
      \begin{bookquote}
        Բավական վատ է, երբ սեփական ուղեղդ է քո դեմ դուրս գալիս։ Իսկ երբ միանում են նաև ուրիշների ուղեղները, դա արդեն ուղղակի ամոթ է։
      \end{bookquote}
      \centering\mascot[1.8cm]{sad}}
    \end{column}
  \end{columns}
\end{frame}
